% Generated draft: requires mathematical and source review.
\documentclass[11pt]{article}
\usepackage[margin=25mm]{geometry}
\usepackage{fontspec}
\setmainfont{Latin Modern Roman}
\usepackage{amsmath,amssymb,mathtools,unicode-math}
\setmathfont{Latin Modern Math}
\usepackage{enumitem,needspace,xurl,hyperref}
\setlength{\parindent}{0pt}
\setlength{\parskip}{5pt}
\setlength{\emergencystretch}{3em}
\newcommand{\N}{\mathbb{N}}
\newcommand{\Z}{\mathbb{Z}}
\newcommand{\Q}{\mathbb{Q}}
\newcommand{\R}{\mathbb{R}}
\newcommand{\C}{\mathbb{C}}
\newcommand{\F}{\mathbb{F}}
\newcommand{\A}{\mathbb{A}}
\newcommand{\PP}{\mathbb{P}}
\newcommand{\E}{\mathbb{E}}
\newcommand{\eps}{\varepsilon}
\newcommand{\dd}{\,\mathrm{d}}
\begin{document}
\section*{UnsolvedMath ID: 30005230}
\textbf{Source problem code:} catalogue entry
\subsection{Definitions and mathematical statement}
For a positive integer \(d\), let \(\Z^d\) denote the integer lattice, with nearest-neighbour edges between vertices whose difference has \(\ell^1\)-norm one. Define its Bernoulli site-percolation threshold by
\[
p_c^{\mathrm{site}}(d)=\inf\bigl\{q\in[0,1]:\mathbb P_q(\text{the origin belongs to an infinite open cluster})>0\bigr\},
\]
where under \(\mathbb P_q\), vertices are independently open with probability \(q\).

The supplied record concerns hostile-environment first-passage percolation (FPPHE) on \(\Z^d\), with type 1 growing from the origin, type 2 seeds of density \(p\), and a type 2 growth parameter \(\lambda\). The excerpt uses a rate-one normalization for comparison with type 1 growth. It does not supply the full stochastic construction: the seed law, passage-time distributions and dependence, activation mechanism, competition rules, and admissible parameter endpoints are absent. Consequently, a uniquely specified probability law cannot be reconstructed from this record alone.

To express the intended question conditionally, fix a completion \(\mathcal M\) of these missing model specifications, let \(\mathbb P^{\mathcal M}_{d,p,\lambda}\) be its law, and let \(I_i(t)\subseteq\Z^d\) be the vertices occupied by type \(i\) at time \(t\). A candidate definition of coexistence is
\[
\mathcal C=\left\{\left|\bigcup_{t\geq0}I_1(t)\right|=\infty\right\}\cap
\left\{\left|\bigcup_{t\geq0}I_2(t)\right|=\infty\right\}.
\]
Here inactive seeds must be distinguished from occupied type 2 vertices if the intended model requires that distinction. The source does not define coexistence or survival, so this candidate must be checked against the original model before acceptance.

\textbf{Requested statement, conditional on those specifications.} Determine whether admissible parameters \(p,\lambda\) exist such that
\[
\mathbb P^{\mathcal M}_{2,p,\lambda}(\mathcal C)>0.
\]
Also determine a coexistence regime satisfying
\[
p<p_c^{\mathrm{site}}(d),\qquad
\mathbb P^{\mathcal M}_{d,p,\lambda}(\mathcal C)>0.
\]
The source does not specify whether the second request is for every \(p\) below the threshold, for some subinterval, or for some parameters, nor whether ``in any dimension'' means each dimension in an intended range or some dimension. These quantifiers remain unresolved; the displayed conditions do not settle them.

\textbf{Supplied contextual qualifications.} The excerpt reports type 1 extinction when \(p>1-p_c^{\mathrm{site}}(d)\), on the stated ground that the origin is almost surely confined to a finite cluster of vertices without type 2 seeds. It also reports extinction for \(\lambda\geq1\) for all \(p\), and more strongly asserts that, for some \(\varepsilon=\varepsilon(d,p)>0\), type 1 growth is contained in first-passage percolation of rate \(1-\varepsilon\), yielding extinction for \(\lambda>1-\varepsilon\). It attributes the comparison to van den Berg and Kesten, reference [1]. These are supplied contextual assertions, not independently verified facts; their seed-density restrictions and endpoint conventions are missing. In particular, the literal phrase ``for all \(p\)'' requires clarification at \(p=0\).

The excerpt says that its Theorem 2 treats small \(p\), that monotonicity is not known, and that survival of type 1 when \(\lambda\) is made small is naturally expected. Its subsequent purported confirmation is truncated. It additionally reports substantially different behaviour on hyperbolic and non-amenable graphs: positive-probability type 2 survival for any \(p\) and \(\lambda\), and an analogue of Theorem 2, citing reference [3]. No graph hypotheses, parameter restrictions, theorem statement, or full reference [3] are supplied. These remarks are background and do not specify the lattice coexistence problem.
\subsection{Short English statement}
The problem asks whether both types can grow indefinitely with positive probability in planar hostile-environment first-passage percolation, and seeks a coexistence regime with seed density \(p<p_c^{\mathrm{site}}(d)\). The supplied record leaves the model construction, the exact coexistence definition, and the second question's quantifiers incomplete. Its open-status label and higher-dimensional results are supplied dated curation, requiring independent review.
\subsection{Sources}
\begin{itemize}
\item Supplied local UnsolvedMath registry record, permanent ID \(30005230\), source problem code \texttt{catalogue entry}, in the Oberwolfach Reports Open Problems collection. No individual proposer is identified; the supplied proposed year is 2022. The record labels the problem open and its statement \texttt{corrected\_verified}; its verification note says only formula delimiters, notation, grammar, and typography were repaired while preserving mathematical content and scope. This is supplied registry metadata, not independent mathematical verification.
\item Supplied report citation: \emph{Large Scale Stochastic Dynamics (2023)}, Oberwolfach Reports, DOI 10.4171/owr/2022/41, pp. 2399--2466. \url{https://doi.org/10.4171/owr/2022/41}. The supplied background describes Oberwolfach Reports as records of mathematical workshops and their open-problem discussions. The citation's 2023 title date and the record's proposed year 2022 are retained as supplied.
\item Finn--Stauffer, hostile-environment first-passage percolation, supplied literature reference: \url{https://ems.press/journals/jems/articles/8736482}.
\item Finn--Stauffer, published PDF, supplied literature reference: \url{https://content.ems.press/assets/public/full-texts/serials/jems/no-issue/8736482/online-first/10.4171-jems-1309-online-first.pdf}.
\item Literature assessment checked at \texttt{2026-08-22}: the supplied curation labels the problem open and reports coexistence proved in dimensions at least three, while the planar case and the full low-density coexistence regime requested here remain unresolved. This assessment is supplied dated curation derived from the OWR record, not independently verified; its time-sensitive status and cited claims require independent checking.
\item The original excerpt attributes a growth comparison to van den Berg and Kesten, reference [1], and statements about hyperbolic and non-amenable graphs to reference [3]. It also invokes Theorem 2. Their complete bibliographic details and theorem statements are not supplied, so these references cannot be resolved from this record alone.
\end{itemize}
\end{document}
